\documentclass[12pt]{article}
\usepackage[T1]{fontenc}
\usepackage[utf8]{inputenc}
\usepackage{lmodern}
\usepackage{textcomp}
\usepackage{enumitem}
\usepackage{geometry}
\geometry{margin=1in}
\begin{document}
\begin{center}
Answer Key - Geography - K-12 Level Assessment 3 \
\small (Answers are ordered exactly as in the question paper.)
\end{center}

\section*{Multiple Choice}
\begin{enumerate}[label=\arabic*.]
\item \textbf{Answer:} C
\\ \textit{Options:} A) structural assimilation.; B) amalgamation theory.; C) acculturation.; D) adaptation.
\\ \textit{Note:} Acculturation refers to the process through which individuals adopt the cultural traits or social patterns of another group, which in this case is the immigrant learning English to integrate into American society.
\item \textbf{Answer:} D
\\ \textit{Options:} A) China.; B) Mexico.; C) Chile.; D) the United States.
\\ \textit{Note:} The United States is the largest exporter of agricultural goods due to its vast and diverse agricultural production capacity, advanced farming technology, and access to international markets.
\item \textbf{Answer:} D
\\ \textit{Options:} A) Dravidian; B) Uralic-Altaic; C) Sino-Tibetan; D) Indo-European
\\ \textit{Note:} The correct answer is Indo-European because it is the language family that encompasses a vast number of languages spoken across North America and Europe, including those from the Baltic, Celtic, Germanic, and Greek branches.
\item \textbf{Answer:} C
\\ \textit{Options:} A) Over 50\% of workforce engaged in agriculture; B) Women held in inferior place in society; C) Main disease related to age and lifestyle; D) High birth and death rates
\\ \textit{Note:} A characteristic of a developed country is that its population tends to experience diseases primarily associated with aging and lifestyle factors, reflecting higher life expectancy and changing health dynamics compared to developing countries.
\item \textbf{Answer:} C
\\ \textit{Options:} A) They have thriving, clean residential areas.; B) There are jobs for most rural-to-urban migrants.; C) They have three separate business districts.; D) They have a well-developed infrastructure.
\\ \textit{Note:} Most African cities are characterized by their distinct economic zones, often comprising three separate business districts that cater to different functions, such as formal, informal, and traditional markets.
\end{enumerate}

\section*{True / False}
\begin{enumerate}[label=\arabic*.]
\setcounter{enumi}{5}
\item \textbf{Answer:} True
\\ \textit{Options:} A) True; B) False
\\ \textit{Note:} The rectangular land survey system, also known as the Public Land Survey System, was developed primarily in the United States to organize land for settlement and is based on a grid system of townships and ranges, differing from the French and Spanish land survey systems which were influenced by their settlement patterns.
\item \textbf{Answer:} True
\\ \textit{Options:} A) True; B) False
\\ \textit{Note:} Abuja, Nigeria, was designated as the capital in 1991 and is situated in the center of the country, but it is not part of the historic central core area, which includes older cities like Lagos and Kano.
\item \textbf{Answer:} True
\\ \textit{Options:} A) True; B) False
\\ \textit{Note:} Boundary problems along the United States-Mexico border are significantly influenced by immigration issues because the movement of people across this border creates tensions related to security, resource allocation, and cultural integration.
\item \textbf{Answer:} True
\\ \textit{Options:} A) True; B) False
\\ \textit{Note:} During the third stage of the demographic transition model, improvements in healthcare, education, and economic conditions lead to a decline in birth rates as families choose to have fewer children, resulting in a slower population growth rate.
\item \textbf{Answer:} True
\\ \textit{Options:} A) True; B) False
\\ \textit{Note:} Double cropping is true because it involves growing two different crops in the same field within a single growing season, allowing farmers to maximize yield and utilize the land more efficiently.
\end{enumerate}

\section*{Long Form Response}
\begin{enumerate}[label=\arabic*.]
\setcounter{enumi}{10}
\item \textbf{Answer:} The core-periphery model of development illustrates how economic resources and opportunities are concentrated in the core regions, while peripheral areas lag behind. Spread effects refer to the positive impacts that occur when growth in the core influences the surrounding peripheral areas, such as job creation, increased investment, and technology transfer. However, increased infrastructure in the core itself is not considered a spread effect because it primarily benefits the core region rather than extending advantages to the periphery. Instead, spread effects arise when the core's development stimulates growth and improvements in the peripheral areas, leading to a more balanced economic landscape.
\\ \textit{Note:} correct because it clearly distinguishes spread effects as the positive influences of core growth on peripheral areas, while emphasizing that improved infrastructure in the core primarily serves to enhance conditions within the core itself, rather than benefiting the surrounding periphery.
\item \textbf{Answer:} Primary economic activities involve the extraction and harvesting of natural resources, such as agriculture, fishing, and mining. Secondary economic activities focus on manufacturing and processing these raw materials into finished goods, while tertiary activities provide services to consumers and businesses. Mining copper is classified as a primary economic activity because it directly involves extracting a natural resource from the earth. This extraction forms the foundational step in the economic process, before any manufacturing or processing occurs. Therefore, since mining involves taking raw materials from the environment, it is categorized as primary rather than secondary.
\\ \textit{Note:} correct because it accurately distinguishes between primary activities, which involve the extraction of natural resources like copper, and secondary activities, which are concerned with manufacturing those resources into finished products.
\end{enumerate}

\vfill
\noindent\textit{End of Answer Key}
\end{document}